\documentclass[10pt,a4paper]{amsart}
\usepackage[margin=19mm]{geometry}
\usepackage{amsmath,amssymb,mathtools}
\usepackage[scr=rsfs,cal=euler]{mathalfa}
\usepackage{xfrac}
\usepackage[unicode,hidelinks,bookmarks=false]{hyperref}
\newcommand{\ie}{i.e.\ }
\newcommand{\cf}{cf.\ }
\newcommand{\resp}{respectively\ }
\newcommand{\Subsection}{\subsection}
\providecommand{\nobreakdash}{\nobreak\hbox{-}}
\setlength{\emergencystretch}{2em}
\multlinegap0pt
\usepackage{amssymb}
\usepackage{mathtools}
\newtheorem{proposition}{Proposition}
\newtheorem{lemma}{Lemma}
\newtheorem{theorem}{Theorem}
\title{Analytic Besov functions, pre-Schwarzian derivatives, \\and integrable Teich\-m\"ul\-ler spaces}
\author{Katsuhiko Matsuzaki, Huaying Wei}
\date{Source: arXiv:2406.13917v5 (CC BY 4.0)}
\begin{document}
\maketitle

\noindent\emph{Source and licence.} Analytic Besov functions, pre-Schwarzian derivatives, \\and integrable Teich\-m\"ul\-ler spaces. Version arXiv:2406.13917v5 (CC BY 4.0), distributed by arXiv under the Creative Commons Attribution 4.0 International licence, http://creativecommons.org/licenses/by/4.0/, as recorded in the arXiv licence metadata for this exact version and shown on the abstract page https://arxiv.org/abs/2406.13917v5. Full licence notice: \texttt{licenses/arxiv-2406.13917-CC-BY-4.0.txt}.\par\smallskip

\noindent\textit{Editorial scope.} Counted unit: the article's theorem bundle (source0.tex lines 1224--1226). The article's complete original proof of it is delivered with it (source0.tex lines 1228--1248, one \texttt{proof} environment of 21 lines). No mathematics is written, repaired or re-derived: the statement and the proof are the article's own lines, in the article's own notation, and no retained line is substituted or re-flowed.  Retained uncounted as the source-local context the proof needs: lines 940--947; lines 1106--1118; lines 1147--1149.  Nothing was omitted from any retained span, so the package carries no omission record.  No retained line contains a citation, so the wrapper carries no bibliography and no bibliography entry.  No retained line contains a figure environment, an image-inclusion command or drawing code, so no image is delivered and the package declares no asset dependency.  Editorial formatting only, in addition: an AMS \texttt{amsart} preamble that restates the article's own declarations and macros used by the retained lines, namely the environments \texttt{proposition}, \texttt{lemma}, \texttt{theorem} on separate counters, together with the standard packages those lines need.  The retained environments print in this document as Proposition 1, Lemma 1, Theorem 1 in order of appearance, whereas the article prints the same items with the numbering of its own full document.  The complete native source of the article as distributed in the arXiv e-print for this exact version is bundled unchanged as \texttt{sources/161055/source0.tex}.
\begin{proposition}\label{S-section}
Let $S:M_p(\mathbb H^+) \to {\mathcal A}_p(\mathbb H^-)$
be the Schwarzian derivative map for $p \geq 1$.
For each $\Psi_0$ in $S(M_p(\mathbb H^+))$, there exists a
neighborhood $V_{\Psi_0}$ of $\Psi_0$ in ${\mathcal A}_p(\mathbb H^-)$ 
and a holomorphic map $\sigma:V_{\Psi_0} \to M_p(\mathbb H^+)$ such that
$S \circ \sigma$ is the identity on $V_{\Psi_0}$.
\end{proposition}

\begin{proposition}\label{affine}
{\rm (i)}
For $\mu, \nu \in M(\mathbb D^*)$, we have $S_{F^{\mu}}=S_{F^{\nu}}$ if and only if  
$F^{\mu}=W \circ F^{\nu}$ on $\mathbb D$
for some M\"obius transformation $W$ of $\widehat{\mathbb C}$ such that  
$W \circ F^{\nu}(\mathbb D)$ is a bounded domain in $\mathbb C$. Moreover, $N_{F^{\mu}}=N_{F^{\nu}}$ if and only if
$F^{\mu}=W \circ F^{\nu}$ on $\mathbb D$
for some affine transformation $W$ of $\mathbb C$.
{\rm (ii)}
For any $\nu \in M_p(\mathbb D^*)$ with $p \geq 1$ and any M\"obius transformation $W$ 
such that $W \circ F^{\nu}(\mathbb D)$ is a bounded domain,  
there exists some $\mu' \in M_p(\mathbb D^*)$ such that $N_{W \circ F^{\nu}}=N_{F^{\mu'}}$. 
\end{proposition}

\begin{lemma}\label{localbound}
$\widetilde L:\widetilde M_p(\mathbb D^*) \to L(M_p(\mathbb D^*))$ is holomorphic.
\end{lemma}

\begin{theorem}\label{bundle}
$J:L(M_p(\mathbb D^*)) \to S(M_p(\mathbb D^*))$ is a holomorphic split submersion for $p \geq 1$. 
\end{theorem}

\begin{proof}
For any $\Phi \in L(M_p(\mathbb D^*))$, let $\Psi_0=J(\Phi) \in S(M_p(\mathbb D^*))$. Then
there exists a
neighborhood $V_{\Psi_0}$ of $\Psi_0$ in $S(M_p(\mathbb D^*))$ 
and a holomorphic map $\sigma:V_{\Psi_0} \to M_p(\mathbb D^*)$ such that
$S \circ \sigma$ is the identity on $V_{\Psi_0}$, as in the case of 
$\mathbb H$ in Proposition \ref{S-section}. Let $\Phi_0=L \circ \sigma(\Psi_0)$, which may be different 
from $\Phi$.
Since $\Phi_0$ can be represented as $\log (F^{\sigma(\Psi_0)})'$, we have
$\Phi=\log (W_a \circ F^{\sigma(\Psi_0)})'$ for some $a \in \widetilde F^{\sigma(\Psi_0)}(\mathbb D^*)$
by Proposition \ref{affine}. 
Namely, $\Phi=\widetilde L(\sigma(\Psi_0),a)$.

Fix this $a$ and define a map
$\widetilde L(\sigma(\cdot),a):V_{\Psi_0} \to L(M_p(\mathbb D^*))$ after shrinking $V_{\Psi_0}$ if necessary.
By Lemma \ref{localbound}, this is a holomorphic map on $V_{\Psi_0}$.
Since $J \circ \widetilde L(\sigma(\Psi),a)=\Psi$ for every $\Psi \in V_{\Psi_0}$, the map
$\widetilde L(\sigma(\cdot),a)$ is a local holomorphic right inverse of $J$ such that $\widetilde L(\sigma(V_{\Psi_0}),a)$ passes through the given 
point $\Phi=\widetilde L(\sigma(\Psi_0),a)$.
This is equivalent to saying that $J$ is a holomorphic split submersion.
\end{proof}

\end{document}
